\documentclass{article}
\usepackage{iclr2024_conference,times}

\usepackage[utf8]{inputenc}
\usepackage[T1]{fontenc}
\usepackage{hyperref}
\usepackage{url}
\usepackage{booktabs}
\usepackage{amsfonts}
\usepackage{nicefrac}
\usepackage{microtype}
\usepackage{titletoc}

\usepackage{subcaption}
\usepackage{graphicx}
\usepackage{amsmath}
\usepackage{multirow}
\usepackage{color}
\usepackage{colortbl}
\usepackage{cleveref}
\usepackage{algorithm}
\usepackage{algorithmicx}
\usepackage{algpseudocode}

\DeclareMathOperator*{\argmin}{arg\,min}
\DeclareMathOperator*{\argmax}{arg\,max}

\graphicspath{{../}}
\begin{filecontents}{references.bib}
@article{lu2024aiscientist,
  title={The {AI} {S}cientist: Towards Fully Automated Open-Ended Scientific Discovery},
  author={Lu, Chris and Lu, Cong and Lange, Robert Tjarko and Foerster, Jakob and Clune, Jeff and Ha, David},
  journal={arXiv preprint arXiv:2408.06292},
  year={2024}
}

@article{wu2016adaptation,
  title={Adaptation in protein fitness landscapes is facilitated by indirect paths},
  author={Wu, Nicholas C and Dai, Lei and Olson, C Anders and Lloyd-Smith, James O and Sun, Ren},
  journal={eLife},
  volume={5},
  pages={e16965},
  year={2016}
}

@article{lin2023evolutionary,
  title={Evolutionary-scale prediction of atomic-level protein structure with a language model},
  author={Lin, Zeming and Akin, Halil and Rao, Roshan and Hie, Brian and Zhu, Zhongkai and Lu, Wenting and Smetanin, Nikita and Verkuil, Robert and Kabeli, Ori and Shmueli, Yaniv and others},
  journal={Science},
  volume={379},
  number={6637},
  pages={1123--1130},
  year={2023}
}

@article{rives2021biological,
  title={Biological structure and function emerge from scaling unsupervised learning to 250 million protein sequences},
  author={Rives, Alexander and Meier, Joshua and Sercu, Tom and Goyal, Siddharth and Lin, Zeming and Liu, Jason and Guo, Demi and Ott, Myle and Zitnick, C Lawrence and Ma, Jerry and others},
  journal={Proceedings of the National Academy of Sciences},
  volume={118},
  number={15},
  year={2021}
}

@article{dallago2021flip,
  title={{FLIP}: Benchmark tasks in fitness landscape inference for proteins},
  author={Dallago, Christian and Mou, Jody and Johnston, Kadina E and Wittmann, Bruce J and Bhattacharya, Nicholas and Goldman, Sam and Bhattacharya, Soumya and Yang, Kevin K},
  journal={bioRxiv},
  year={2021}
}

@article{stanton2022accelerating,
  title={Accelerating Bayesian optimization for biological sequence design with denoising autoencoders},
  author={Stanton, Samuel and Maddox, Wesley and Gruver, Nate and Maffettone, Phillip and Delaney, Emily and Greenside, Peyton and Wilson, Andrew Gordon},
  journal={arXiv preprint arXiv:2203.12742},
  year={2022}
}

@article{notin2022tranception,
  title={Tranception: protein fitness prediction with autoregressive transformers and inference-time retrieval},
  author={Notin, Pascal and Dias, Mafalda and Fraternali, Franca and Marchetti, Jonathan and Listgarten, Jennifer and Gal, Yarin},
  journal={arXiv preprint arXiv:2205.13760},
  year={2022}
}

@article{meier2021language,
  title={Language models enable zero-shot prediction of the effects of mutations on protein function},
  author={Meier, Joshua and Rao, Roshan and Verkuil, Robert and Liu, Jason and Sercu, Tom and Rives, Alexander},
  journal={Advances in Neural Information Processing Systems},
  volume={34},
  pages={29287--29303},
  year={2021}
}
\end{filecontents}

\title{Protein Fitness Prediction on the GB1 Landscape Using Frozen ESM-2 Embeddings}

\begin{document}

\maketitle

\begin{abstract}
Predicting protein fitness from sequence is a central challenge in protein engineering and directed evolution. We study the GB1 combinatorial fitness landscape—a benchmark dataset of ~150,000 four-residue mutants with experimentally measured fitness scores—using frozen ESM-2 protein language model embeddings as input features to a shallow MLP. By decoupling representation learning (handled by ESM-2, pre-trained on evolutionary data) from fitness regression, we obtain a computationally efficient and interpretable baseline. We report Spearman correlation and MSE on held-out sequences and use this framework to investigate how properties of the embedding space, training set composition, and model architecture affect generalization across the fitness landscape.
\end{abstract}

\section{Introduction}
\label{sec:intro}

Protein engineering relies on efficiently navigating fitness landscapes—mappings from amino acid sequence to biological function. Combinatorial mutagenesis experiments such as the GB1 dataset of \citet{wu2016adaptation} provide dense measurements over a small sequence space and have become standard benchmarks for evaluating machine learning models \citep{dallago2021flip}.

Protein language models (PLMs) trained on large sequence databases have demonstrated strong transfer to downstream prediction tasks \citep{rives2021biological, lin2023evolutionary}. ESM-2 \citep{lin2023evolutionary} in particular produces rich per-residue representations that encode evolutionary, structural, and physicochemical information. However, the degree to which these representations capture fitness—a property not directly observed during pre-training—remains an open question.

In this work we establish a clean baseline: extract mean-pooled ESM-2 embeddings (frozen weights) and train a small MLP to predict log-fitness. This separates the contribution of the PLM representation from that of supervised training and provides a platform for controlled experiments on the GB1 landscape.

\section{Background}
\label{sec:background}

\subsection{GB1 Fitness Landscape}
GB1 is an IgG-binding domain whose fitness landscape was measured by \citet{wu2016adaptation} via deep mutational scanning of four residue positions (V39, D40, G41, V54). The resulting dataset contains all single, double, triple, and quadruple mutants (20$^4$ = 160,000 possible variants; ~149,000 observed), with fitness measured as enrichment relative to wild type.

\subsection{ESM-2 Protein Language Model}
ESM-2 \citep{lin2023evolutionary} is a family of masked language models trained on 250M UniRef90 sequences at scales from 8M to 15B parameters. Its representations have been shown to encode secondary structure, contact maps, and mutational effects \citep{meier2021language}.

\section{Methods}
\label{sec:methods}

\subsection{Dataset and Preprocessing}
We use the public GB1 dataset from \citet{wu2016adaptation}. Fitness values are log-transformed ($\log(1 + f)$) to reduce skew. We split 80\% for training and 20\% for evaluation, stratified randomly with three seeds for reliability estimates.

\subsection{Frozen ESM-2 Embeddings}
For each variant sequence, we pass the full protein sequence through ESM-2 with frozen weights and extract the final-layer representations. We mean-pool over the four mutated positions to obtain a fixed-dimensional embedding vector. No gradient flows through ESM-2 during training.

\subsection{MLP Architecture and Training}
A two-layer MLP (hidden dimensions 256 and 128, ReLU activations, dropout 0.1) maps embeddings to scalar fitness predictions. We minimize MSE loss with Adam (lr=1e-3, weight decay 1e-4) and a cosine annealing scheduler over 50 epochs.

\subsection{Evaluation Metrics}
We report MSE and Spearman rank correlation ($\rho$) on held-out sequences. Spearman $\rho$ is the primary metric as it captures ranking ability, which is most relevant for directed evolution applications.

\section{Experimental Results}
\label{sec:results}

\begin{figure}[h]
    \centering
    \begin{subfigure}{0.48\textwidth}
        \includegraphics[width=\textwidth]{train_loss.png}
        \caption{Training loss (MSE) over epochs.}
    \end{subfigure}
    \hfill
    \begin{subfigure}{0.48\textwidth}
        \includegraphics[width=\textwidth]{val_spearman.png}
        \caption{Validation Spearman correlation over epochs.}
    \end{subfigure}
    \caption{Learning curves for the ESM-2 + MLP baseline on the GB1 fitness landscape, averaged over 3 random seeds. Shaded regions indicate standard error.}
    \label{fig:learning_curves}
\end{figure}

\begin{figure}[h]
    \centering
    \includegraphics[width=0.5\textwidth]{final_spearman_bar.png}
    \caption{Final held-out Spearman correlation ($\pm$ stderr over 3 seeds).}
    \label{fig:final_spearman}
\end{figure}

\section{Conclusion}
\label{sec:conclusion}

We present a clean, reproducible baseline for protein fitness prediction on the GB1 landscape using frozen ESM-2 embeddings and a shallow MLP. This framework decouples representation learning from supervised regression and enables controlled investigation of what drives fitness predictability. Future work can explore ESM-2 model scale, fine-tuning strategies, training set composition effects, and the geometry of the fitness landscape in embedding space.

\bibliography{references}
\bibliographystyle{iclr2024_conference}

\end{document}
